% !TEX program = xelatex
\documentclass[12pt,a4paper]{article}
\usepackage{fontspec}
\usepackage{polyglossia}
\setmainlanguage{english}
\setmainfont{DejaVu Serif}
\setsansfont{DejaVu Sans}
\setmonofont{DejaVu Sans Mono}
\usepackage{geometry}
\geometry{margin=2.5cm}
\usepackage{amsmath,amssymb,amsthm,mathtools,bm}
\usepackage{microtype}
\usepackage{enumitem}
\usepackage{hyperref}
\hypersetup{colorlinks=true, linkcolor=blue, citecolor=blue, urlcolor=blue}
\newtheorem{definition}{Definition}
\newtheorem{axiom}{Axiom}
\newtheorem{theorem}{Theorem}
\newtheorem{proposition}{Proposition}
\newtheorem{remark}{Remark}
\DeclareMathOperator{\Valid}{ValidInterface}
\DeclareMathOperator{\Interface}{Interface}
\DeclareMathOperator{\Auto}{Auto}
\DeclareMathOperator{\AlwaysAct}{AlwaysAct}
\DeclareMathOperator{\Input}{Input}
\DeclareMathOperator{\Adm}{Adm}
\DeclareMathOperator{\Dir}{Dir}
\DeclareMathOperator{\AutoAct}{AutoAct}
\DeclareMathOperator{\PassiveChannel}{PassiveChannel}
\DeclareMathOperator{\Control}{Control}
\DeclareMathOperator{\Der}{Der}
\DeclareMathOperator{\Struct}{Struct}
\DeclareMathOperator{\Dist}{Dist}
\DeclareMathOperator{\End}{End}
\DeclareMathOperator{\Tr}{Tr}
\DeclareMathOperator{\Attr}{Attr}
\DeclareMathOperator{\Coh}{Coh}
\DeclareMathOperator{\Sync}{Sync}
\DeclareMathOperator{\Loss}{Loss}
\DeclareMathOperator{\Stab}{Stab}
\newcommand{\M}{\mathcal M}
\newcommand{\B}{\mathcal B}
\newcommand{\Hcal}{\mathcal H}
\newcommand{\Ispace}{\mathfrak I}
\newcommand{\XiAuto}{\Xi_I^{\mathrm{auto}}}
\newcommand{\XiMan}{\mathcal M_{\Xi}^{\mathrm{auto}}}
\newcommand{\NXi}{\mathcal N_{\Xi}}
\newcommand{\PiAmp}{\Pi_{\mathrm{amp}}}
\newcommand{\PiStr}{\Pi_{\mathrm{str}}}
\newcommand{\PiStruct}{\Pi_{\mathrm{struct}}}

\title{\textbf{TMI-Core 0.2: The Axiomatic Core of Automatic Interfaces}\\
\large The Primary Interface, Definitional Automaticity, and the Transition to Automatically Structured Fields}
\author{Salkutsan Aleksey}
\date{2026}

\begin{document}
\maketitle

\begin{abstract}
This document fixes the updated core of the Theory of Manifold Interfaces, denoted by \(\mathcal T_{MI}^{0.2}\). The central transition of version 0.2 is that automaticity is no longer treated as an external addition to an interface. It is introduced as a definitional property of a valid interface: if a structure cannot autonomously interpret an admissible input and produce an intelligible output for another system, then it is not a full interface, but a passive channel or a component of a meta-interface. The core connects the primary interface axiom, definitional automaticity, the distinction between an interface and a passive channel, derived interfaces, projection into the physical layer, and the scaling of automatically structured interface fields.
\end{abstract}

\tableofcontents
\newpage

\section{Purpose of the Document}
The purpose of this document is to fix a minimal core that should not be expanded by adding new major entities. Further work should refine definitions, prove internal consequences, and derive testable physical or engineering projections.

\section{Primary Objects}
\begin{definition}[Null Undifferentiation]
\(N_0\) is the limiting state without distinction, sides, relation, order, or direction. It is not a physical vacuum and not an object of nature, but a limiting symbol of undifferentiation inside the axiomatic scheme.
\end{definition}

\begin{definition}[Primary Interface]
\(I_0\) is the first boundary of differentiation through which a minimal distinction arises.
\end{definition}

\begin{definition}[First Distinction]
\(\Delta_0\) is the minimal result of the action of the primary interface. It unfolds into the minimal two-sided structure \((A,\neg A)\).
\end{definition}

\section{The Primary Interface Axiom}
\begin{axiom}[Primary Interface]
\begin{equation}
\boxed{A_{I_0}:\quad N_0\xrightarrow{I_0}\Delta_0\to(A,\neg A).}
\end{equation}
A shorter form is
\begin{equation}
\boxed{N_0\xrightarrow{I_0}(A,\neg A).}
\end{equation}
\end{axiom}

The meaning of the axiom is that the primary interface does not connect already existing sides; it creates the very possibility of sides. Before \(I_0\), there is no \(A\) and no \(\neg A\). After the boundary arises, the distinction \(\Delta_0\) appears, and this distinction unfolds into the minimal structure \((A,\neg A)\).

\section{Valid Interface}
\begin{definition}[Valid Interface]
A valid interface \(I\) is a structure that, in the presence of an admissible input, performs translation, interpretation, or structuring between systems without an external controller directing every act of translation.
\end{definition}

In working form:
\begin{equation}
I_{A\to B}\equiv\bigl(\partial(A,B),T_I(S_A\to S_B),\operatorname{Interp}_I,\Auto_I\bigr),
\end{equation}
where \(\partial(A,B)\) is the boundary of interaction, \(T_I\) is the state-translation operator, \(\operatorname{Interp}_I\) is the operation of interpretation, and \(\Auto_I\) is the self-execution of the interface function under an admissible input.

\section{The Axiom of Interface Automaticity}
\begin{axiom}[Definitional Automaticity]
\begin{equation}
\boxed{A_{\mathrm{auto}}:\quad \forall I\in\Ispace,\;\Valid(I)\Rightarrow\Auto(I).}
\end{equation}
In short:
\begin{equation}
\boxed{I\in\Ispace\Rightarrow\Auto(I).}
\end{equation}
\end{axiom}

Automaticity does not mean permanent activity:
\begin{equation}
\boxed{\Auto(I)\not\Rightarrow\AlwaysAct(I).}
\end{equation}

Automaticity means self-activation when an admissible input is present:
\begin{equation}
\boxed{\Auto(I)\equiv\bigl(\Input(I)\land\Adm(I)\land\Dir(I)\Rightarrow\AutoAct(I)\bigr).}
\end{equation}

\section{Passive Channel and Meta-Interface}
\begin{definition}[Passive Channel]
A passive channel is a structure through which a signal or a state can pass, but which does not perform its own interpretive function without an external controlling agent.
\end{definition}

If translation is controlled by an external agent \(U_{\mathrm{ext}}\), then the valid interface is the composite structure
\begin{equation}
\boxed{I'=(I,U_{\mathrm{ext}}).}
\end{equation}

\begin{equation}
\boxed{\PassiveChannel(I)\land\Control(U_{\mathrm{ext}},I)\Rightarrow\Interface(I').}
\end{equation}

\begin{theorem}[Interface-Channel Distinction]
If a structure does not perform the interface function automatically under an admissible input, then it is either not a valid interface or it is a passive channel included in a broader meta-interface:
\begin{equation}
\boxed{\neg\Auto(I)\Rightarrow \neg\Valid(I)\lor\PassiveChannel(I).}
\end{equation}
\end{theorem}

\begin{proof}
By the axiom of automaticity, a valid interface satisfies \(\Valid(I)\Rightarrow\Auto(I)\). Its contrapositive gives \(\neg\Auto(I)\Rightarrow\neg\Valid(I)\). If the structure nevertheless participates in transmission but does not perform translation by itself, then it is a passive channel inside a broader system \(I'=(I,U_{\mathrm{ext}})\). Therefore the statement holds inside the axiomatic scheme.
\end{proof}

\section{Derived Interfaces}
\begin{definition}[Derived Interface]
\(\Der(I,I_0)\) means that the interface \(I\) is derived from the primary interface \(I_0\), i.e. it belongs to the space of interfaces that becomes possible after the primary differentiation.
\end{definition}

\begin{equation}
\boxed{\Der(I,I_0)\Rightarrow I\in\Ispace\Rightarrow\Auto(I).}
\end{equation}

For derived interfaces, automaticity is unfolded conditionally:
\begin{equation}
\boxed{\Der(I,I_0)\Rightarrow\bigl(\Input(I)\land\Adm(I)\land\Dir(I)\Rightarrow\AutoAct(I)\bigr).}
\end{equation}

\section{Bridge to the Physical Layer}
If a physical field is an interface field,
\begin{equation}
\Xi_I\in\Ispace_{\mathrm{phys}},
\end{equation}
then it is automatic as an interface:
\begin{equation}
\boxed{\Xi_I\in\Ispace_{\mathrm{phys}}\Rightarrow\Auto(\Xi_I).}
\end{equation}
Consequently, the physical interface field should be considered as automatically structured:
\begin{equation}
\boxed{\Xi_I\in\Ispace_{\mathrm{phys}}\Rightarrow\XiAuto.}
\end{equation}

\section{Projection Chain}
The full form of the field is
\begin{equation}
\boxed{\XiAuto.}
\end{equation}
Its projections are
\begin{equation}
\boxed{\XiAuto\to\Xi_I\to\chi_I.}
\end{equation}
Here \(\chi_I\) is the minimal scalar amplitude projection:
\begin{equation}
\boxed{\chi_I=\PiAmp(\XiAuto).}
\end{equation}

\section{Scaling}
A local automatically structured field scales into a network and a manifold:
\begin{equation}
\boxed{\XiAuto\to\{\Xi_{I,k}^{\mathrm{auto}}\}_{k\in K}\to\NXi\to\XiMan\to\Struct_I(\M).}
\end{equation}

\section{Unified Core Formula}
\begin{equation}
\boxed{N_0\xrightarrow{I_0}\Delta_0\to(A,\neg A)\Rightarrow\Ispace\Rightarrow\Auto(I).}
\end{equation}
\begin{equation}
\boxed{I_0\Rightarrow\Ispace\Rightarrow\Auto(I)\Rightarrow\XiAuto\Rightarrow\XiMan\Rightarrow\Struct_I(\M).}
\end{equation}

\section{Status of the Claims}
\begin{itemize}
\item Level A: axioms and definitions, for example \(\Valid(I)\Rightarrow\Auto(I)\).
\item Level B: internal theorems, for example the transition from \(\XiAuto\) to \(\Struct_I\) under stabilization conditions.
\item Level C: physical hypotheses, for example the existence of \(\XiAuto\) as a physical mechanism.
\item Level D: testable consequences, for example an increase of coherence time in a quantum experiment.
\end{itemize}

\section{Conclusion}
TMI-Core 0.2 fixes the central spine of the theory: the primary interface creates the possibility of distinction; a valid interface is automatic by definition; physical interface fields are automatically structured; and a collection of such fields forms a scalable interface manifold capable of stabilizing as structure.

\begin{thebibliography}{9}
\bibitem{TMIcore} Salkutsan Aleksey. \textit{Theory of Manifold Interfaces. An Axiomatic Conceptual-Mathematical Framework of Interfaciality, Meaning, and Cognition}. Zenodo. \href{https://doi.org/10.5281/zenodo.19940675}{doi:10.5281/zenodo.19940675}.
\bibitem{TMIphysics} Salkutsan Aleksey. \textit{TMI-Physics 2.2: Automatically Structured Interface Fields}. 2026.
\end{thebibliography}
\end{document}
